\chapter{背景と目的}
%該当分野の従来の状況，問題点，本プロジェクトで設定した課題，実施した解決策，及び成果を簡潔に記述する．


\section{プロジェクトの立ち上げ}

世の中にはニーズを十分に満たしていないシステムが存在する．これは，開発者が作るものとユーザが求めるものとの乖離が原因の1つであると考えられる．この問題を解決するためには，ユーザを理解してシステムを開発する必要があることから，開発者が現場に赴き，調査をして，ユーザから直接学ぶべきであると考えた．そこで「使ってもらって学ぶフィールド指向システムデザイン」を理念とするプロジェクトが始まった．
\bunseki{大須賀雅也}

\section{プロジェクトの方針}
本プロジェクトは，例年実際に現場に赴くフィールド調査とアジャイル開発手法の1つである，スクラム手法\cite{scrumb}を採用している．フィールド調査ではユーザの思考や行動などの，現場に行かないと分からないことを知ることができる．また，スクラムは，アジャイル開発の中でも少人数のチームで開発を行い，スプリントと呼ばれる固定の短い時間に区切って作業を進めるものである\cite{scrumb}．これはユーザのフィードバックを繰り返し受けて，改善する機会を何度も得ることができるため，本プロジェクトの「使ってもらって学ぶフィールド指向システムデザイン」という理念と合致する．以上のことから，今年度もフィールド調査とスクラム手法を採用することにした．
\bunseki{大須賀雅也}

\section{未来大生支援グループ}\label{sec:gaiyou}
% 上述の問題点を解決すべく当プロジェクトの掲げる課題の概要を述べる．
私達未来大生支援グループは初め，函館市の病院に関連するアプリを開発することを目的として結成された．病院に関連するアプリについてアイディアを出す過程で，課題に対する具体的なアイディアが出尽くしたと感じ，このままでは興味深い結果が得られないと考えるようになった．この状況を受け，一度全てを白紙に戻し，全体を再度検討することとした．そして，私達自身に有用なアプリの開発に焦点を当てることを決定した．私達全員が日常的に利用し，かつ改善の余地があると感じるものを対象に検討した結果，私たち自身が在籍する公立はこだて未来大学に注目することとなった．普段から大学のシステムを頻繁に利用している私達は，その操作性の悪さに共通認識を持っていた．この認識を背景に，再度ブレインストーミングを行った．その結果，大学のシステムを単に改善するだけでなく，大学生活の利便性を高めることを目指すことになった．具体的には，私達が抱える問題を解決するため，公立はこだて未来大学の情報を一元化したアプリを開発することを決定した．

\bunseki{大須賀雅也}
